% 05-mandatory-courses.tex -- Mandatory Scientific Training
\chapter{Mandatory Scientific Training}
\label{ch:mandatory-courses}

The six mandatory courses form the scientific backbone of the program.  Each
course carries four ECTS credits, for a total of twenty-four ECTS.  They are
designed to establish a common methodological foundation so that all doctoral
candidates---regardless of their disciplinary background---can engage
confidently with the research challenges of digital finance.

\vspace{6pt}

\coursecard{MST-01}{Foundation of Data Science}{4}{BBU}{%
  Students will acquire the core skills needed to work with large-scale
  financial datasets.  The course covers data wrangling, exploratory analysis,
  statistical modeling, and reproducible research workflows in Python and~R.
  Participants learn to source and clean financial data from public and
  proprietary databases, apply descriptive and inferential statistics, and
  produce publication-quality visualizations.  Emphasis is placed on
  reproducibility through version control and literate programming.
  \\[4pt]
  \textbf{Prerequisites:} Basic statistics, working knowledge of Python or~R.\\
  \textbf{Assessment:} Project-based.
}

\vspace{4pt}

\coursecard{MST-02}{Introduction to AI for Financial Applications}{4}{WWU}{%
  Participants learn to design, train, and evaluate machine learning models for
  financial applications.  Topics include supervised and unsupervised learning,
  feature engineering for financial data, cross-validation strategies that
  respect temporal structure, and the fundamentals of model deployment.
  Students will implement end-to-end ML pipelines, from data preprocessing
  through model selection to backtesting on historical market data.
  \\[4pt]
  \textbf{Prerequisites:} Data Science Foundations or equivalent.\\
  \textbf{Assessment:} Coding assignments and a final project.
}

\vspace{4pt}

\coursecard{MST-03}{The Need for eXplainable AI: Methods and Applications in Finance}{4}{BFH}{%
  Students will master techniques for making AI-driven financial decisions
  transparent and auditable.  The course covers post-hoc explainability
  methods---SHAP, LIME, counterfactual explanations---as well as inherently
  interpretable models.  Participants learn to quantify fairness using
  established metrics, assess model bias in lending and insurance contexts,
  and map their findings to the compliance requirements of the EU~AI~Act.
  \\[4pt]
  \textbf{Prerequisites:} AI for Finance or equivalent.\\
  \textbf{Assessment:} XAI audit report on a real-world financial model.
}

\vspace{4pt}

\coursecard{MST-04}{Introduction to Blockchain Applications in Finance}{4}{ASE}{%
  Participants explore distributed ledger technologies and their impact on
  financial services.  The course introduces DLT fundamentals, consensus
  mechanisms, and smart-contract programming before examining decentralized
  finance (DeFi) protocols, asset tokenization, and central bank digital
  currencies (CBDCs).  Students will develop and deploy a smart contract on a
  test network and conduct a comparative regulatory analysis of European
  blockchain frameworks.
  \\[4pt]
  \textbf{Prerequisites:} None beyond program admission.\\
  \textbf{Assessment:} Smart-contract project and regulatory analysis paper.
}

\vspace{4pt}

\coursecard{MST-05}{Sustainable Finance}{4}{UNA}{%
  Students will develop a thorough understanding of the tools, frameworks, and
  regulations that connect financial markets to sustainability objectives.
  Topics include ESG rating methodologies, climate-risk scenario analysis, green
  bonds and sustainability-linked instruments, the EU~Taxonomy, and the
  Sustainable Finance Disclosure Regulation~(SFDR).  Participants learn to
  evaluate corporate sustainability reports and to integrate ESG factors into
  investment decision-making.
  \\[4pt]
  \textbf{Prerequisites:} None beyond program admission.\\
  \textbf{Assessment:} Sustainability assessment report.
}

\vspace{4pt}

\coursecard{MST-06}{Ethics Applicable to Digital Aspects}{4}{UTW}{%
  Students examine the ethical dimensions of artificial intelligence and
  digital technologies in financial services.  Topics include algorithmic
  fairness, accountability frameworks, transparency requirements, and human
  oversight principles as they apply to AI-driven financial decision-making.
  Participants learn to conduct ethical impact assessments and map their
  analyses to the EU~AI~Act's risk classification system.
  \\[4pt]
  \textbf{Prerequisites:} None beyond program admission.\\
  \textbf{Assessment:} Ethical impact assessment report.
}
